\section{Первопричина и главная теорема}\label{sec:firstcause}

Всякая редукция из \cref{sec:reduction} давит внутрь, к единственному открытому узлу
$\code{TheLastStep00Obligation}$, к тому же суженному до масштаба $A \ge 5$. Дальше внутрь идти
некуда. Настоящий раздел разворачивается и смотрит наружу, отвечая на три вопроса. Почему последний
узел законно принять как внешнюю первопричину, а не как скрытую теорему? Что такая архитектура
доказывает \emph{безусловно}, без всякой аксиомы? И, самое глубокое, \emph{почему} узел нельзя
закрыть изнутри --- почему знание его причины недостижимо для наблюдателя, живущего внутри системы?
Ответ на третий вопрос и есть наша главная теорема.

\subsection{Причину нельзя вывести изнутри}

С открытым узлом можно поступить двояко: оставить открытым (честно, но рассказ обрывается) или
принять (и тогда всё, что от него зависело, замыкается). Второй путь законен ровно при одном условии
--- принятие должно быть \emph{видимым}: можно объявить явное допущение и проследить, что на нём
держится, но нельзя тайком выдать недоказанное за доказанное. Оказывается, свободы в том, \emph{как}
принять, нет. Принять узел изнутри --- предъявить его внутреннее самообоснование --- невозможно, и
это теорема.

Назовём внутреннее доказательство, пересекающее собственную границу, чтобы обосновать её же,
$\code{InternalSelfDerivationOfStep00CausalClosure}$. Такая конструкция строит запрещённый конкретный
евклидов двигатель, $\code{InternalSelfDerivationOfStep00CausalClosure} \to
\code{SomeConcreteEuclideanEngine}$; но двигателей нет (ацикличность ранга $\code{lexRank}$,
\cref{sec:engine}), и потому её не существует.

\begin{theorem}[\code{no\_internalSelfDerivation\_step00CausalClosure}]\label{thm:no_internalSelfDerivation_step00CausalClosure}
\green{} Причинная граница не самозаверяется:
$\lnot\,\code{InternalSelfDerivationOfStep00CausalClosure}$.
\end{theorem}

\emph{Идея доказательства.} Самозаверение отождествляется, шаг за шагом машинной проверкой, с
построением бесконечно убывающей чистой цепи --- ровно евклидова двигателя. Двигателей нет, значит нет
и самозаверения.

То же на языке начала: самопричинённый первый кадр, интернализированное событие $0 \to 1$, не может
существовать в стабильной безмоторной архитектуре.

\begin{theorem}[\code{no\_internalisedHorizonBoundary}]\label{thm:no_internalisedHorizonBoundary}
\green{} Интернализированной причинной границы горизонта не существует:
$\lnot\,\code{InternalisedStep00HorizonBoundary}$.
\end{theorem}

\begin{remark}
Сам по себе запрет самозаверения близок к тавтологии: доказать $P$, пересекая границу $P$,
сворачивается в $P \wedge \lnot P$. Содержательно здесь \emph{отождествление формы} --- что попытка
самообоснования оказывается уже сожжённой двигательной конструкцией, а не какой-то новой
невозможностью. Физически это вечный двигатель первого рода: самообоснование извлекало бы работу
вывода из ничего. Ранг $\code{lexRank}$ --- сохраняющаяся величина со строгим дном, и петля обрывается
об это дно. Поэтому фундамент нужно принести снаружи.
\end{remark}

\subsection{Аксиома, принятая намеренно: событие \texorpdfstring{$0\to1$}{0->1}}

Раз изнутри причину положить нельзя, мы кладём её извне, намеренно, как единственную аксиому всего
построения, \axiom{}. Её тип $\code{Step00FirstCause}$ --- структура события $0 \to 1$, перехода от
небытия к первому причинному кадру. Два из пяти её полей --- чистые маркеры, несущие лишь
$\mathrm{True}$: $\code{origin}$ (сингулярность $0$, где внутреннего языка ещё нет) и
$\code{firstFrame}$ (первый кадр $1$). Остальные три --- содержательные причинные границы, по одной на
каждую ветвь программы:

\begin{definition}[три причинные границы \axiom{}]\label{def:step00FirstCause}
$\code{Step00FirstCause}$ несёт $\code{causalBoundary}$ --- границу близнецов, в точности открытый
узел $\code{TheStrictLastStep00Obligation}$; $\code{riemannBoundary}$ --- закон Римана о манифестации
нуля (\cref{sec:riemann}); и $\code{nsBoundary}$ --- гейт-закон Навье--Стокса энергобаланса
(\cref{sec:ns}). Вся математическая тяжесть живёт в этих трёх полях; маркеры не утверждают ничего.
\end{definition}

Оформление меняет происхождение результата, но не его силу: декрет --- ровно сумма того, что в него
положено.

\begin{theorem}[\code{step00FirstCause\_iff\_causalClosure}]\label{thm:step00FirstCause_iff_causalClosure}
\green{} Первопричина эквивалентна конъюнкции своих границ: $\code{Step00FirstCause}
\leftrightarrow (\text{узел} \wedge \text{закон Римана} \wedge \text{гейт НС})$.
\end{theorem}

Аксиома заперта в \emph{карантинном модуле}, и карантин --- машинная процедура: отдельный верификатор
помечает каждую зависящую от \axiom{} декларацию как \code{AXIOM-TAINTED}, и ни одна не протекла в
зелёную линию. Прежняя аксиома замыкания теперь теорема, добываемая проектированием,
$\code{step00CausalClosure} := \axiom{}.\code{causalBoundary}$. По близнецовой ветви она даёт
бесконечность нижних twin-центров --- \emph{условно}.

\begin{theorem}[\code{twinLowersInfinite\_from\_step00CausalClosure}]\label{thm:twinLowersInfinite_from_step00CausalClosure}
\yellow{} При принятой границе $\code{TwinLowers.Infinite}$.
\end{theorem}

\emph{Это условный вывод, а не доказательство гипотезы близнецов.} Более того, декрет не слабее своего
следствия: он предъявляет twin-центр выше любого порога,
$\forall M_0,\ \exists m,\ M_0 < m \wedge \code{TwinCenterZ}\,m$
(\code{causalClosureAxiom\_asserts\_twins\_at\_every\_scale}, \yellow{}). Принять эту аксиому
\emph{значит} принять близнецов; никакой бесплатной силы она не создаёт. Безаксиомный остаток честен
до тавтологии.

\begin{theorem}[\code{nonAxiomaticRemainingObligation\_iff\_lastStep00Obligation}]\label{thm:nonAxiomaticRemainingObligation_iff_lastStep00Obligation}
\green{} Безаксиомный остаточный долг равен старому узлу:
$\code{Step00NonAxiomaticRemainingObligation} \leftrightarrow \code{TheLastStep00Obligation}$
\red{}. Никакая обёртка не опускает открытое содержание ниже узла;
$\code{twin\_prime\_conjecture}$ остаётся $\code{sorry}$ и через карантин не замыкается.
\end{theorem}

\subsection{Эпистемика: причина есть, узнать нельзя, знание всё ломает}

Обозначим через $\code{InternalKnowledgeOfCause}$ внутреннее выведение причины. О ней можно сказать
три вещи. Она есть --- это \axiom{}, \yellow{}, принятая намеренно. Её нельзя узнать --- и это
теорема, \green{}, а не оговорка.

\begin{theorem}[\code{cause\_unknowable}]\label{thm:cause_unknowable}
\green{} Внутреннее знание причины невозможно:
$\lnot\,\code{InternalKnowledgeOfCause}$.
\end{theorem}

\emph{Идея доказательства.} Знать причину изнутри значило бы вывести её изнутри, а это строит вечный
двигатель ($\code{knowledge\_builds\_perpetualEngine}$:
$\code{InternalKnowledgeOfCause} \to \code{SomeConcreteEuclideanEngine}$), которого нет.
Непознаваемость причины --- не слабость, а закон сохранения.

Знание, будь оно возможно, было бы взрывным: из него следует $\lnot\,\code{TwinLowers.Infinite}$
($\code{knowledge\_finitizes\_twins}$) \emph{и} $\code{TwinLowers.Infinite}$
($\code{knowledge\_proves\_anything}$), оба \green{}, оба \emph{ex falso} через невозможный двигатель.
Содержательный, невзрывной остаток --- дихотомия
$\lnot\,\code{InternalKnowledgeOfCause} \vee \lnot\,\code{TwinLowers.Infinite}$
($\code{unknowable\_or\_twins\_finite}$, \green{}), чей левый дизъюнкт есть
\cref{thm:cause_unknowable}: мы всегда в левом мире.

Несущее содержание --- не взрыв, а аксиомо-свободная лемма, сердце причинной линии, где гипотеза
«двигателей нет» потребляется содержательно, а не вакуумно.

\begin{theorem}[\code{twins\_infinite\_of\_noEngine\_and\_boundary}]\label{thm:twins_infinite_of_noEngine_and_boundary}
\green{} Отсутствие двигателей вместе с причинной границей влечёт бесконечность близнецов:
\begin{equation}\label{eq:twins_infinite_of_noEngine}
\lnot\,\code{SomeConcreteEuclideanEngine} \to \code{Step00CausalClosureAxiom} \to
\code{TwinLowers.Infinite}.
\end{equation}
\end{theorem}

\begin{proof}
Предположим, близнецы конечны; тогда есть последняя граница $M_0$, выше которой их нет. Из этой
конечной границы строится бесконечная семья порождённых потоков; конечный ключ вынужден по пигеонхолу
дать коллизию; из коллизии собирается конкретный двигатель-свидетель. Именно его убивает гипотеза
$\lnot\,\code{SomeConcreteEuclideanEngine}$ --- не пустой взрыв, а прямое попадание. Гипотеза
потреблена содержательно.
\end{proof}

Инстанциируя \cref{eq:twins_infinite_of_noEngine} поставляемым $\code{lexRank}$ отсутствием двигателей
и границей из аксиомы, получаем близнецов декретом.

\begin{theorem}[\code{twins\_because\_unknowable}]\label{thm:twins_because_unknowable}
\yellow{} $\code{TwinLowers.Infinite}$, причём вывод видимо проходит через ту же «двигателей нет»,
что даёт непознаваемость.
\end{theorem}

\begin{remark}
«Близнецы бесконечны, потому что причину узнать нельзя» означает \emph{через общую причину и несущую
лемму}, а не «непознаваемость в одиночку доказывает бесконечность». Вся тяжесть \emph{существования}
по-прежнему лежит на принятой границе: непознаваемость без границы близнецов не даёт --- иначе они были
бы уже доказаны.
\end{remark}

\subsection{Две стены}

Есть и вторая, независимая причина, по которой близнецы ускользают от наблюдателя --- познавательная,
а не причинная, и безусловная (ядро опирается лишь на $\code{propext}$). Зафиксируем корректный
сертификат $\code{FiniteSystemKnowsTwin}$, зависящий только от конечного вида уровня $A$. Тогда
конечное знание --- это знание про класс: конечная система знает близнеца $B$, только если весь его
класс сито-эквивалентности состоит из близнецов ($\code{knowledge\_forces\_pure\_class}$, \green{});
близнец в смешанном классе принципиально невидим ($\code{mixed\_class\_twin\_unknowable}$, \green{}); а
при хвостово смешанных классах бесконечность вообще не удостоверить
($\code{infinitude\_unknowable\_of\_eventually\_mixed}$, \green{}), с тривиальным видом как самым
резким безусловным случаем ($\code{trivialView\_infinitude\_unknowable}$, \green{}). То, что классы
\emph{нетривиального} вида на конкретных уровнях действительно смешиваются, --- арифметический вход
\red{}; структурные теоремы и близорукая инстанциация безусловны. Обе стены замыкаются в одно
утверждение.

\begin{theorem}[\code{two\_walls\_one\_nature}]\label{thm:two_walls_one_nature}
\green{} Изнутри конечного вида не видно близнеца со смешанным классом, изнутри системы не видно
первопричины:
\[
\begin{aligned}
&\bigl(\forall S,A,\mathrm{Cert},B,\mathrm{bad}:\ \code{SieveEquivalent}\,S\,A\,B\,\mathrm{bad} \to \lnot\,\code{IsTwin}\,\mathrm{bad}\\
&\qquad {}\to \lnot\,\code{FiniteSystemKnowsTwin}\,S\,A\,\mathrm{Cert}\,B\bigr)\ \wedge \lnot\,\code{InternalKnowledgeOfCause}.
\end{aligned}
\]
\end{theorem}

Обе стены --- одно и то же: конечная система не может выйти за собственную границу. Конечный вид уровня
$A$ --- познавательный горизонт, как строгий порядок обхода --- причинный.

\subsection{Главная теорема: высшая энергетическая несовместимость}

Всё собирается в одно утверждение, склеивающее причинную стену с познавательной, целиком в зелёном
ядре.

\begin{theorem}[\code{higherEnergyIncompatibility\_main}]\label{thm:higherEnergyIncompatibility_main}
\green{} (без аксиом) Пять граней одной несовместимости выполняются совместно:
(1)~знание первопричины изнутри строит вечный двигатель;
(2)~поэтому первопричина непознаваема изнутри;
(3)~конечный вид знает близнеца, только если весь его класс --- близнецы;
(4)~при хвостовом смешении классов бесконечность изнутри неудостоверима;
(5)~несущая грань --- сама эта несовместимость (двигателей нет $=$ причину не узнать) вместе с
принятой границей влечёт $\code{TwinLowers.Infinite}$.
\end{theorem}

Грани (1)--(4) --- это две стены бок о бок; грань (5) --- мост от них к выводу. Прочитать
энергетически: знать изнутри стоит вечного двигателя, которого нет, а бесконечность близнецов ---
внешнее знание, оплаченное первопричиной. Инстанциируя грань (5) декретом, получаем следствие ---
жёлтое, над зелёным ядром.

\begin{corollary}[\code{higherEnergyIncompatibility\_twins}]\label{thm:higherEnergyIncompatibility_twins}
\yellow{} $\code{TwinLowers.Infinite}$. \emph{Это не доказательство гипотезы близнецов};
$\code{twin\_prime\_conjecture}$ остаётся $\code{sorry}$, не замыкающимся через карантин.
\end{corollary}

\begin{remark}
Название не украшение. Принцип Ландауэра делает информацию физичной --- стереть бит стоит не меньше
$kT\ln 2$, а демон Максвелла соблюдает второй закон именно потому, что его знание приходится
оплачивать. Наша теорема --- дискретный собрат: узнать причину изнутри стоило бы бесконечной работы
(вечного двигателя), значит, запрещено тем же законом сохранения; бесконечность близнецов не
бесплатна, её цена --- вынесенное наружу знание, то есть аксиома. Отсюда три мира, которые
расчерчивает декрет. Если узел истинен, аксиома лишняя, а близнецы на деле доказаны; если узел
независим, теория непротиворечива, но узнать это изнутри нельзя; если узел опровержим, карантин
противоречив, и растяжка ($\code{quarantine\_inconsistent\_if\_node\_refuted}$, \yellow{}) срабатывает
разом, обесценивая всякую заражённую декларацию. Оба внутренних разрешения стоят вечного двигателя ---
потому границу можно принять только извне.
\end{remark}
